\documentclass[11pt,a4paper]{article}
\usepackage[italian]{babel}
\usepackage[margin=2.5cm]{geometry}
\usepackage{parskip}
\usepackage{amsmath, amssymb, amsthm}
\usepackage[most]{tcolorbox}
\usepackage{hyperref}

% --- Riquadri con bordo nero, senza numero ---
\tcbset{nota/.style={colback=white, colframe=black, boxrule=0.5pt,
  sharp corners}}
\NewTColorBox{definizione}{}{nota,
  before upper={\textbf{Definizione.}\ }}
\NewTColorBox{teorema}{o}{nota,
  before upper={\textbf{Teorema\IfValueT{#1}{ (#1)}.}\ }}
\NewTColorBox{proposizione}{o}{nota,
  before upper={\textbf{Proposizione\IfValueT{#1}{ (#1)}.}\ }}

% --- Ambienti semplici (senza riquadro) ---
\theoremstyle{definition}
\newtheorem*{esempio}{Esempio}
\newtheorem*{osservazione}{Osservazione}

\title{Congruenza modulo $n$}
\author{Marco Cerbella}
\date{Lezione 3 -- 28 settembre 2026}

\begin{document}
\maketitle

\section{La congruenza modulo $n$}

\begin{definizione}
Siano $a, b \in \mathbb{Z}$. Si dice che $a$ \emph{divide} $b$, e si scrive
$a \mid b$, se
\[
\exists\, k \in \mathbb{Z} \text{ tale che } b = ka.
\]
Se non esiste un tale $k$ si scrive $a \nmid b$ (``$a$ non divide $b$'').
\end{definizione}

Sia $n \geq 1$ un intero fissato.

\begin{definizione}
La \emph{congruenza modulo $n$} è la relazione $\equiv_n$ su $\mathbb{Z}$
definita da
\[
a \equiv_n b \iff n \mid (b - a) \iff \exists\, k \in \mathbb{Z} \text{
tale che } b - a = kn \quad (b = kn + a).
\]
Si scrive anche $a \equiv b \pmod n$.
\end{definizione}

\begin{proposizione}
$\equiv_n$ è una relazione di equivalenza su $\mathbb{Z}$, e
\[
\mathbb{Z}/_{\equiv_n} = \mathbb{Z}_n = \{[0]_n, [1]_n, \dots, [n-1]_n\}.
\]
\end{proposizione}

\begin{proof}
Verifichiamo che $\equiv_n$ è una relazione di equivalenza.
\begin{enumerate}
    \item \emph{Riflessiva}: $a \equiv_n a$, perché $n \mid (a - a) = 0$,
          cioè $0 = k \cdot n$ con $k = 0$.
    \item \emph{Simmetrica}: se $a \equiv_n b$, allora $n \mid (b - a)$,
          cioè esiste $k \in \mathbb{Z}$ con $b - a = kn$. Posto
          $k' = -k$ si ha $a - b = -kn = k'n$, quindi $n \mid (a - b)$,
          cioè $b \equiv_n a$.
    \item \emph{Transitiva}: siano $a \equiv_n b$ e $b \equiv_n c$.
          Allora $b - a = kn$ e $c - b = hn$ per certi $k, h \in
          \mathbb{Z}$. Sommando,
          \[
          c - a = kn + hn = (k + h)n,
          \]
          quindi $n \mid (c - a)$, cioè $a \equiv_n c$. \qedhere
\end{enumerate}
\end{proof}

La classe di $a$ è
\begin{align*}
[a]_n &= \{b \in \mathbb{Z} \mid a \equiv_n b\}
       = \{b \in \mathbb{Z} \mid n \mid (b - a)\} \\
      &= \{b \in \mathbb{Z} \mid \exists\, k \in \mathbb{Z} \text{ tale
         che } b = kn + a\}
       = \{kn + a \mid k \in \mathbb{Z}\},
\end{align*}
cioè è l'insieme di tutti i numeri che differiscono da $a$ per un multiplo
di $n$. L'insieme quotiente è
$\mathbb{Z}_n = \{[0]_n, [1]_n, \dots, [n-1]_n\}$.

\begin{esempio}
Per $n = 3$ le classi $[0]_3, [1]_3, [2]_3$ sono gli insiemi $X_0, X_1,
X_2$ dell'esempio delle partizioni di $\mathbb{Z}$ (lezione 2): $\mathbb{Z}_3$ è proprio la partizione
di $\mathbb{Z}$ in multipli di $3$, multipli di $3$ più $1$ e multipli di
$3$ più $2$.
\end{esempio}

\end{document}
